%Dokumentklasse

%draft als optionohne bilder für bessere performance
%\documentclass[a4paper,12pt,]{scrreprt}

%normal mit Bildern
\documentclass[
a4paper,
12pt,
draft=True]
{scrartcl}

%Section als Chapter
\RedeclareSectionCommand[%
%beforeskip = -1sp plus -1sp minus -1sp,% kleinster negativer Wert, um den Absatzeinzug nach der Überschrift zu verhindern.
afterskip = 1.5 \baselineskip plus -1sp minus 1sp,
font = \Huge,
]{section}

\usepackage[left= 3cm,right = 3cm, bottom = 3cm,top = 3cm]{geometry}
%\usepackage[onehalfspacing]{setspace}

% ============= Packages =============
% Dokumentinformationen
\usepackage[
pdftitle={Praktikum - Umwelttechnik},
pdfsubject={},
pdfauthor={Roman-Luca Zank},
pdfkeywords={},	
%Links nicht einrahmen
hidelinks
]{hyperref}

%nur Text zum prüfen des Umfangs

% Standard Packages
%\usepackage[bottom]{footmisc}
\usepackage[utf8]{inputenc}
\usepackage[ngerman]{babel}

\usepackage[T1]{fontenc}
%\usepackage{helvet}

%\renewcommand{\familydefault}{\sfdefault}

\usepackage{graphicx}
\graphicspath{{img/}}
\usepackage{mhchem}
\usepackage{fancyhdr}
\usepackage{lmodern}
\usepackage{color}
\usepackage{placeins}
\usepackage{booktabs}
\usepackage{caption}
\usepackage[list=true]{subcaption}
\usepackage{longtable}
\usepackage{tikz}
\usepackage{pgfplots}
\usepackage{lastpage}
%\usepackage{ulem}
\usepackage{mathtools}
\usepackage{adjustbox}
\usetikzlibrary{patterns}
\usepackage{pdfpages}

%Einheitenpackage
\usepackage{siunitx}  
\sisetup{	locale = DE, 
	per-mode=fraction,
	inter-unit-product=\ensuremath{\cdot},
	detect-weight = true,
	quotient-mode=fraction
}
%neue Einheiten definieren
\DeclareSIUnit\xyz{xyz}	
\DeclareSIUnit\rpm{rpm}	
\DeclareSIUnit\mws{mWS}	
\DeclareSIUnit\degrees{^\circ}	

%Automatisch cdot statt *
\DeclareMathSymbol{*}{\mathbin}{symbols}{"01}


%Tabelle
\usepackage{tabularx}
\usepackage{tabulary}

%nur letzte Zeile der Gleichung nummerieren
\makeatletter
\def\Let@{\def\\{\notag\math@cr}}
\makeatother

% zusätzliche Schriftzeichen der American Mathematical Society
\usepackage{amsfonts}
\usepackage{amsmath}

%Abkürzungsverzeichnis
\usepackage{acronym}

%kein Abstand bei neuem Kapitel vom Seitenanfang
%\vspace*{2.3\baselineskip} = ORIGINAL
%\renewcommand*{\chapterheadstartvskip}{\vspace*{.0\baselineskip}}

%nicht einrücken nach Absatz
\setlength{\parindent}{0pt}


\urlstyle{same}

% ============= Kopf- und Fußzeile =============
\pagestyle{fancy}
%
\lhead{}
\chead{}
\rhead{}%\slshape }%\leftmark}
%%
\lfoot{}
\cfoot{}
\rfoot[{\thepage\ of \pageref*{LastPage}}]{Seite \thepage\ von \pageref*{LastPage}}
%%
\renewcommand{\headrulewidth}{0pt}
\renewcommand{\footrulewidth}{0pt}
%\renewcommand{\chapterpagestyle}{fancy}

%Fußnotelinie
%\let\footnoterule

%Fußnote mit Klammer
\renewcommand*{\thefootnote}{(\arabic{footnote})}

%Abb. statt Abbildung
\addto\captionsngerman{%
	\renewcommand{\figurename}{Abb.}%
	\renewcommand{\tablename}{Tab.}%
}

% ============= Package Einstellungen & Sonstiges ============= 
%Besondere Trennungen
%\hyphenation{De-zi-mal-tren-nung}
\usepackage[none]{hyphenat}
\hyphenpenalty=5000
\tolerance=5000
\providecommand\phantomsection{}

\usepackage{mathtools}


% ============= Dokumentbeginn =============

\begin{document}
%Seiten ohne Kopf- und Fußzeile sowie Seitenzahl
\pagestyle{empty}

%\begin{center}
\begin{tabular}{p{\textwidth}}


\begin{center}
\includegraphics[scale=0.75]{logos.jpg}\\
\end{center}


\\

\begin{center}
\LARGE{\textsc{
Protokoll \\
Strömungslehre\\
}}
\end{center}

\\

%\begin{center}
%\large{Fakultät für Muster und Beispiele \\
%der Hochschule Musterhausen \\}
%\end{center}
%
%\\

\begin{center}
\textbf{\Large{Versuch Rohrwand}}
\end{center}

\begin{center}
	\large{Gruppe II (BCUT3)}
\end{center}


\\
%\begin{center}
%zur Erlangung des akademischen Grades\\
%Bachelor of Engineering
%\end{center}


%\begin{center}
%vorgelegt von
%\end{center}

\begin{center}
\Large{\textbf{Teilnehmer:}} \\ 
\end{center}
\begin{center}
\large{Carolin Lohmann \\
	Willy Messerschmidt \\
	Marcus Willbrandt \\
	Roman-Luca Zank} \\
\end{center}


\\

\begin{center}
\begin{tabular}{lll}
\large{\textbf{Protokollführer:}} & & \large{Carolin Lohmann \&} \\
&& \large{Roman-Luca Zank}\\
&&\\
\large{\textbf{Datum der Versuchsdurchführung:}}&& \large{16.12.2019}\\
&&\\
\large{\textbf{Abgabedatum:}}&& \large{13.01.2020}
\end{tabular}
\end{center}

\\ \\ \\ \\ \\ 
\large{Merseburg den \today}

\end{tabular}
\end{center}


%\include{14_danksagungen}

%\include{15_zusammenfassung}

% Beendet eine Seite und erzwingt auf den nachfolgenden Seiten die Ausgabe aller Gleitobjekte (z.B. Abbildungen), die bislang definiert, aber noch nicht ausgegeben wurden. Dieser Befehl fügt, falls nötig, eine leere Seite ein, sodaß die nächste Seite nach den Gleitobjekten eine ungerade Seitennummer hat. 
\cleardoubleoddpage

% Pagestyle für Titelblatt leer
\pagestyle{empty}

%Seite zählen ab
\setcounter{page}{0}

%Titelblatt
\begin{center}
\begin{tabular}{p{\textwidth}}


\begin{center}
\includegraphics[scale=0.75]{logos.jpg}\\
\end{center}


\\

\begin{center}
\LARGE{\textsc{
Protokoll \\
Strömungslehre\\
}}
\end{center}

\\

%\begin{center}
%\large{Fakultät für Muster und Beispiele \\
%der Hochschule Musterhausen \\}
%\end{center}
%
%\\

\begin{center}
\textbf{\Large{Versuch Rohrwand}}
\end{center}

\begin{center}
	\large{Gruppe II (BCUT3)}
\end{center}


\\
%\begin{center}
%zur Erlangung des akademischen Grades\\
%Bachelor of Engineering
%\end{center}


%\begin{center}
%vorgelegt von
%\end{center}

\begin{center}
\Large{\textbf{Teilnehmer:}} \\ 
\end{center}
\begin{center}
\large{Carolin Lohmann \\
	Willy Messerschmidt \\
	Marcus Willbrandt \\
	Roman-Luca Zank} \\
\end{center}


\\

\begin{center}
\begin{tabular}{lll}
\large{\textbf{Protokollführer:}} & & \large{Carolin Lohmann \&} \\
&& \large{Roman-Luca Zank}\\
&&\\
\large{\textbf{Datum der Versuchsdurchführung:}}&& \large{16.12.2019}\\
&&\\
\large{\textbf{Abgabedatum:}}&& \large{13.01.2020}
\end{tabular}
\end{center}

\\ \\ \\ \\ \\ 
\large{Merseburg den \today}

\end{tabular}
\end{center}
 %Prokolle
%\begin{center}
\begin{tabular}{p{\textwidth}}


\begin{center}
\includegraphics[scale=0.75]{img/logos.jpg}\\
\end{center}


\\

\begin{center}
\LARGE{\textsc{
Recherche \\
Rückgewinnung von Ammoniak aus Industrieabwässern\\
}}
\end{center}

%\begin{center}
%\large{Fakultät für Muster und Beispiele \\
%der Hochschule Musterhausen \\}
%\end{center}
%
%\\
 \\
 
\begin{center}
\textbf{\Large{Seminararbeit in Medienrecherche}}
\end{center}

\begin{center}
	\large{im WiSe 2019}
\end{center}
 \\
%\begin{center}
%zur Erlangung des akademischen Grades\\
%Bachelor of Engineering
%\end{center}


\begin{center}
\large{vorgelegt von}
\end{center}
\\


\begin{center}
\Large{\textbf{Roman-Luca Zank}} \\
\end{center}

\begin{center}
3. Semester \\
Chemie- und Umwelttechnik \\
\end{center}


\begin{center}
\begin{tabular}{lll}
	\textbf{E-Mail:} & & romanzank@mail.de\\
	\textbf{Matrikelnummer:} & &25240\\
	\textbf{Adresse:} & &Platz der Bausoldaten 2, Zimmer 224\\
	\textbf{Ort:} & &06217 Merseburg\\
	&& \\
	\textbf{Prüfer:} & & Dr. Frank  Baumann\\
\end{tabular}
\end{center}

\\ \\ \\ \\ \\
\large{Merseburg, \today}

\end{tabular}
\end{center}
 %Seminar-/Abschlussarbeit

% Pagestyle für Rest des Dokuments
\pagestyle{fancy}

%Inhaltsverzeichnis
\tableofcontents
\thispagestyle{empty}

%Inhalt
%
%Verzeichnis aller Bilder
\label{sec:bilder}
\listoffigures
\addcontentsline{toc}{chapter}{Abbildungsverzeichnis}
\thispagestyle{empty}

%Verzeichnis aller Tabellen
\label{sec:tabellen}
\listoftables
\addcontentsline{toc}{chapter}{Tabellenverzeichnis}
\thispagestyle{empty}



%Abkürzungsverzeichnis
%\setlength{\columnsep}{20pt}
%\twocolumn
%\addchap{Nomenklatur}
%\label{sec:abkurzung}
%\begin{acronym}
%\acro{kf}[$\text{k}_\text{f}$]{Durchlässigkeitsbeiwert}
%\acro{t}{Durchlaufzeit}
%\acro{tm}[$\text{t}_\text{m}$]{Mittlere Durchlaufzeit}
%\acro{V}{Volumen}
%\acro{h}{Höhe der Wassersäule}
%\acro{Q}{Volumenstrom}
%\acro{l}{Durchströmte Länge}
%\acro{A}{Grundfläche}
%\acro{d}{Durchmesser}
%
%\end{acronym}
%\subsubsection{Aufrufen einer Abkürzung}
%\acs{rT}
%\begin{verbatim}
%\acs{Abkürzung}
%\end{verbatim}

\section{Zielbeschreibung}
\label{sec:zielbeschreibung}

\subsection{Kinetik}
Die kinetischen Parameter einer Verseifungsreaktion sollen aus der Messgröße Leitfähigkeit bestimmt werden.
Die betrachteten Größen sind Reaktionsordnung, Stoßfaktor und Aktivierungsenergie.
\subsection{Verweilzeit}
Ziel dieses Versuches ist die Beschreibung einer realen Reaktorkaskade. Aus experimentellen Daten sollen die Kenngrößen Bodensteinzahl und Kesselzahl des Verweilzeitverhaltens abgeleitet werden.


\subsection{Umsatzgrad}
Der Umsatzgrad der Verseifungsreaktion, welche auch in den vorausgegangenen Versuchen betrachtet wurde, wird in diesem Versuch an einer realen Anlage unter identischen Prozessbedingungen wie auch zuvor ermittelt.

\section{Physikalische Hintergründe}
\label{sec:physik}

Als physikalische Hintergründe sind die folgenden Gleichungen dargestellt. Diese beschränken sich im Wesentlichen auf die \textsc{Bernoulli}-Gleichung mit deren Ableitungen, die \textsc{Reynold}szahl, sowie dem K$_V$-Wert.\\

\textsc{Bernoulli}-Gleichung:
\begin{flalign}
	p_1+z_1*g*\rho +\frac{1}{2}*\rho*(v_1)^2&= p_2+z_2*g*\rho+\frac{1}{2}*\rho*(v_2)^2
\end{flalign}

\textsc{Bernoulli}-Gleichung mit Druckverlust $\Delta p_v$:
\begin{flalign}
p_1+z_1*g*\rho +\frac{1}{2}*\rho*(v_1)^2&= p_2+z_2*g*\rho+\frac{1}{2}*\rho*(v_2)^2\boldsymbol{+\Delta p_v}
\end{flalign}

Vereinfachung mit $p_{\text{geodätisch}} = const.$ und $p_{\text{dynamisch}} = const.$:
\begin{flalign}
p_1&= p_2+\Delta p_v
\end{flalign}

Kontinuitätsgleichung:
\begin{flalign}
	\dot{V}	&= v*A\\
	\dot{V_1}&= \dot{V_2}\\
	v_1*A_1	&= v_2*A_2
\end{flalign}

Druckverluste:
\begin{flalign}
	\Delta p_v	&= \left(\lambda *\frac{l}{d}+\sum\zeta_i\right)*\frac{\rho}{2}*v^2
\end{flalign}

\textsc{Reynold}szahl:
\begin{flalign}
	Re	&= \frac{v*d*\rho}{\eta} = \frac{v*d}{\nu}
\end{flalign}

K$_V$-Wert:
\begin{flalign}
K_V	&= \dot{V}*\sqrt{\frac{\SI{1}{\bar}}{\Delta p}*\frac{\rho}{\rho_0}}
\end{flalign}

\section{Geräte und Chemikalien}
\label{sec:geraete}

\textbf{Geräte:}
\begin{itemize}
	\item Magnetrührer mit Rührfisch
	\item Bechergläser
	\item Erlenmeyerkolben
	\item Büchnertrichter
	\item Saugflasche
	\item Filterpapier
	\item Reflektometer RQflex$^{\textsuperscript{\textregistered}}$ plus 10 von \textsc{Merck}
\end{itemize}

\vspace*{5mm}

\textbf{Proben/Chemikalien:}
\begin{itemize}
	\item destilliertes Wasser
	\item Abwasserproben 1, 2 \& 3
		\item Schnelltests von \textsc{Chemsolute$^{\textsuperscript{\textregistered}}$}:
	\begin{itemize}
		\item Phosphat Test für \SI{0}{} - \SI{500}{\milli \gram \per \liter} \ce{PO4^3-} (Art.-Nr. 29500001) 
		\item Nitrat Test für \SI{0}{} - \SI{500}{\milli \gram \per \liter} \ce{NO3-} (Art.-Nr. 29350001) 
		\item Nitrit Test für \SI{0}{} - \SI{80}{\milli \gram \per \liter} \ce{NO2-} (Art.-Nr. 29300001) 
	\end{itemize}
	\item Reflectoquanten$^{\textsuperscript{\textregistered}}$ von \textsc{Merck} für Reflektometer:
	\begin{itemize}
		\item Ammonium Test für \SI{0.2}{} - \SI{7.0}{\milli \gram \per \liter} \ce{NH4+} (Art.-Nr. 1168920001) 
		\item Phosphate Test für \SI{5}{} - \SI{120}{\milli \gram \per \liter} \ce{PO4^3-} (Art.-Nr. 1169780001) 
		\item Nitrat Test für \SI{3}{} - \SI{90.0}{\milli \gram \per \liter} \ce{NO3-} (Art.-Nr. 1169950001) 
		\item Nitrit Test für \SI{0.5}{} - \SI{25.0}{\milli \gram \per \liter} \ce{NO2-} (Art.-Nr. 1169730001) 
	\end{itemize}
\end{itemize}







\section{Durchführung}
\label{sec:durchfuerung}

\subsection{Durchführung zum Versuch Kinetik}
Der Verwendete Reaktor wies ein Volumen von \SI{1}{\liter} auf. Das Reaktionsgemisch wurde mit einem Ankerrührer bei einer Drehzahl von \SI{100}{\per\minute} homogenisiert. Ein Thermostat sorgte für konstante Temperaturen. 
Der Versuch wurde drei mal wiederholt. Dabei lagen die Temperaturen bei \SI{23,5}{\degreeCelsius}, \SI{30,2}{\degreeCelsius} und \SI{37,0}{\degreeCelsius}.

Es wurde die hydrolytische Spaltung des Essigsäureethylesters in alkalischem Millieu untersucht. 

Einer Vorlage aus \SI{1}{\liter} 0,05 molarer Natriumhydroxidlösung wurde zum Zeitpunkt Null jeweils 0,05 Mol reiner Essigsäureethylether zugegeben.   

Die Leitfähigkeitsmessung erfolgte automatisch alle 5 Sekunden, durch das \textsc{inoLab-IDS Multi\nolinebreak 9430}, in Verbindung mit einer entsprechenden Sonde, automatisch, für annähernd 15 Minuten.

\subsection{Durchführung zum Versuch Verweilzeit}

Anstelle des einfachen Rührkessels wurde hier eine Rührkesselkaskade aus drei Rührkesseln eingesetzt. Jeder Kessel verfügte über einen eigenen Rührer. Es wurde deionisiertes Wasser und eine eingefärbte, schwache Salzlösung verwendet. 
Im ersten Teilversuch wurde durch die Zugabe einer geringen Menge Salzlösung eine Impulsmarkierung erzeugt.  
Der Volumenstrom an deionisiertem Wasser belief sich auf \SI{8}{\liter\per\hour}.\linebreak
Im zweiten Teil fand eine Verdrängung des deionisierten Wassers durch das leicht salzbeladene Wasser statt. Der Salzwasservolumenstrom betrug in diesem Fall \SI{10}{\liter\per\hour}.

Messstellen für die Leitfähigkeit befanden sich hinter jedem Kessel. \\Die Leitfähigkeitswerte wurden automatisch alle fünf Sekunden über den Computer erfasst.
Zur Dosierung des Wassers wurde über Schlauchpumpen verwirklicht.
\subsection{Durchführung zum Versuch Umsatzgrad}

In Vorbereitung auf den Versuch zum Umsatzgrad wurden je \SI{6}{\liter} 0,1 molare Natriumhydroxid- und Essigsäureethylesterlösung angemischt. Die Anlage mit der Rührkesselkaskade wurde entleert. Es wurde der gleiche Gesamtvolumenstrom wie im zweiten Teil des Versuches zur Verweilzeit, von \SI{10}{\liter\per\hour}, eingestellt. Dieser setzte sich aus \SI{5}{\liter\per\hour} Natriumhydroxidlösung und \SI{5}{\liter\per\hour} Essigsäuereethylesterlösung zusammen. Die Leitfähigkeitswerte wurden nach jedem der drei Kessel automatisch alle fünf Sekunden aufgenommen.






\section{Auswertung}
\label{sec:auswertung}

%Darstellung Messungen und Kurzbeschreibung

\begin{figure}[h!]
	\begin{center}
		%\resizebox{0.8\textwidth}{!}{
		\begin{tikzpicture}[trim axis left, trim axis right]
		\begin{axis}[
		axis lines = left,
		width = 13cm,
		height = 7cm,
		xmin = 0,
		xmax = 20,
		ymin = 0,
		ymax = 0.1,
		%ytick = {0,2,...,14},
		%xtick = {0,10,...,100},
		ylabel={c in [\si{\mole\per\liter}]},
		xlabel={t in [\si{\second}]},
	%	legend style={at={(0.5,0.95)},anchor=west}
		]
		\addplot table {rt-M-x-replace-string.dat};
		\end{axis}
		\end{tikzpicture}
		%	}Ventilkennlinie
		\caption{Konzentration über der Zeit}
		\label{dia:c/t}
	\end{center}
\end{figure}
\FloatBarrier


\section{Fazit}
\label{sec:fazit}




\newpage
\section{Fehlerbetrachtung}
\label{sec:fehler}

Bei der Herstellung der Natronlauge für den Versuchsteil Kinetik ging etwas Lösung verloren. Dies wurde dadurch ersichtlich, dass beim befüllen von drei Maßkolben zu je \SI{1}{\litre}, ein Maßkolben nicht gänzlich gefüllt werden konnte. Der Verlust wird auf maximal \SI{5}{\milli\liter} geschätzt. Die resultierende Abweichung beträgt damit 0,5\% und kann gegenüber anderen zufälligen Fehlern vernachlässigt werden. 

Im Versuchsteil Umsatzgrad wurde der Verlauf dadurch gestört, dass eine Schlauchpumpe Luft zog. Die Reaktionslösung setzte sich damit ab einem schwer zu bestimmenden Zeitpunkt nicht mehr äquimolar zusammen. Die Endleitfähigkeit konnte daher nicht aus dem Produktgemisch ermittelt werde. Anstatt dieser wurde für die Berechnungen die mittlere Endleitfähigkeit aus dem Versuchsteil Kinetik noch einmal verwendet. 

Der Essigsäureethylether besitzt einen Flammpunkt von \SI{-4}{\degreeCelsius} und ist damit als leicht flüchtig einzustufen. Während der Umsetzung im offenen Rührkessel, insbesondere bei Temperaturen bis knapp \SI{40}{\degreeCelsius}, könnte ein geringer Anteil verdampft sein. Da die zugegebene Menge auf $\frac{1}{100}$\si{\milli\liter} genau abgemessen wurde, könnte sich dieser Einfluss auf die Reaktion ausgewirkt haben.

Messfehler könnten ihren Ursprung in den Sensoren haben. Die Verwendung von mindestens 4 verschiedenen Leitfähigkeitsmessgeräten, welche vorher nicht mit einer einheitlichen Kalibrierlösung überprüft wurden, ist als bedeutende Fehlerquelle anzunehmen.  

Während der Messungen im Rührkessel schwankte die Temperatur trotz Thermostat um rund $\pm$\SI{0,5}{\degreeCelsius}. 
Ein exaktes Ablesen der eingestellten Volumenströme an der Rührkesselkaskade war nicht möglich, da die Durchflussanzeigen, durch die stoßweise Förderung der Schlauchpumpen, nie zur Ruhe kamen.

Die Reaktion wurde im Rührkessel jeweils für nur etwa 15 Minuten beobachtet. Die Reaktion war zu diesem Zeitpunkt noch nicht vollständig zum erliegen gekommen. Die Messung der Endleitfähigkeit hätte noch später stattfinden müssen, da die Reaktion möglicherweise noch, wenn auch nur noch sehr wenig, lief.
%Die anschließende Lagerung der Reaktionslösungen, vor der Messung der Endleitfähigkeit, bei Raumtemperatur verfälscht das Messergebnis der Endleitfähigkeit möglicherweise, da die auf \SI{30}{\degreeCelsius} und \SI{37}{\degreeCelsius} geheizten Proben dabei abkühlten. 

\section*{Anhang}
\addcontentsline{toc}{chapter}{Anhang}
\label{sec:anhang}

%Tabelle START
\vspace*{-2.5mm}
\renewcommand{\arraystretch}{1.2}
\begin{table}[h!]
	\centering
	\caption{Dichte des Wassers zu unterschiedlichen Temperaturen mittels \cite{BernhardSpang.2002}}
	\label{tab:dichte}
	%\resizebox{10cm}{!}{
	\begin{tabulary}{\textwidth}{C|C}
		\hline
		\textbf{Temperatur} & \textbf{Dichte mittels \texttt{=densW(T,p)}} \\ 
		\hline
		\SI{25,4}{\celsius} & \SI{996,98}{\kg\per\raiseto{3}\meter} \\
		\SI{25,5}{\celsius} & \SI{996,96}{\kg\per\raiseto{3}\meter} \\
		\SI{25,6}{\celsius} & \SI{996,93}{\kg\per\raiseto{3}\meter} \\
		\SI{25,7}{\celsius} & \SI{996,91}{\kg\per\raiseto{3}\meter} \\
		\SI{25,8}{\celsius} & \SI{996,88}{\kg\per\raiseto{3}\meter} \\
		\SI{25,9}{\celsius} & \SI{996,85}{\kg\per\raiseto{3}\meter} \\
		\SI{26,0}{\celsius} & \SI{996,83}{\kg\per\raiseto{3}\meter} \\
		\SI{26,1}{\celsius} & \SI{996,80}{\kg\per\raiseto{3}\meter} \\
		\SI{26,2}{\celsius} & \SI{996,77}{\kg\per\raiseto{3}\meter} \\
		\SI{26,3}{\celsius} & \SI{996,75}{\kg\per\raiseto{3}\meter} \\
		\SI{26,4}{\celsius} & \SI{996,72}{\kg\per\raiseto{3}\meter} \\
		\SI{26,5}{\celsius} & \SI{996,69}{\kg\per\raiseto{3}\meter} \\
		\SI{26,6}{\celsius} & \SI{996,67}{\kg\per\raiseto{3}\meter} \\
		\SI{26,7}{\celsius} & \SI{996,64}{\kg\per\raiseto{3}\meter} \\
		\SI{26,8}{\celsius} & \SI{996,61}{\kg\per\raiseto{3}\meter} \\
		\SI{26,9}{\celsius} & \SI{996,58}{\kg\per\raiseto{3}\meter} \\
		\SI{27,0}{\celsius} & \SI{996,56}{\kg\per\raiseto{3}\meter} \\
		\SI{27,1}{\celsius} & \SI{996,53}{\kg\per\raiseto{3}\meter} \\
		\SI{27,2}{\celsius} & \SI{996,50}{\kg\per\raiseto{3}\meter} \\
		\SI{27,3}{\celsius} & \SI{996,48}{\kg\per\raiseto{3}\meter} \\
		\SI{27,4}{\celsius} & \SI{996,45}{\kg\per\raiseto{3}\meter} \\
		\SI{27,5}{\celsius} & \SI{996,42}{\kg\per\raiseto{3}\meter} \\
		\hline
	\end{tabulary}
	%}
\end{table}
\FloatBarrier
\vspace*{-2.5mm}
%Tabelle Ende

%Praktikumsskript, Modul ………, Versuch …….., Prof. Musterprof. 
%DIN 12345, Jahr der Veröffentlichung 
%Link der Internetseite, Zugriffsdatum 
%Buchtitel, Autor, Verlag, Veröffentlichungsjahr 

%Literaturverzeichnis Bücher
\bibliography{Literatur}
\bibliographystyle{unsrtdin}




%\chapter*{Eidesstattliche Erklärung}
\label{erklaerung}
Hiermit versichere ich, die vorliegende Seminararbeit selbstständig und nur unter Verwendung der von mir angegebenen Quellen und Hilfsmittel verfasst zu haben. Sowohl inhaltlich als auch wörtlich entnommene Inhalte wurden als solche kenntlich gemacht. Die Arbeit hat in dieser oder vergleichbarer Form noch keinem anderem Prüfungsgremium vorgelegen. \\
\\[1.5cm]
Datum:	\hrulefill\enspace Unterschrift: \hrulefill
\\[3.5cm]
\addcontentsline{toc}{chapter}{Selbstständigkeitserklärung}

\end{document}
